% CellCharter研究专题结论与展望模块（非整篇硕士学位论文总论终稿）
% 对应真实E1/E2,E3c,E3d,E3b,E5a,E5b结果；不预设新算法或统计显著性。
\section{本研究的结论与展望}
\label{sec:cc-conclusion-module-v1}

\subsection{主要实证结论}
本研究在成熟CellCharter空间邻域表示与聚类流程基础上，围绕特征深度、聚类数确定、空间边界以及跨组织评价开展了系统的计算实验。DLPFC部分包含12张切片、3名独立供体；外部数据为6名乳腺癌患者。本研究保存各阶段冻结实验方案、技术种子、预测来源哈希、完整正负结果和事后分析记录，避免仅凭最有利切片或单次初始化得出过度概括的结论。

第一，在固定Oracle-K条件下，将DLPFC每张切片的技术种子先平均再对12切片等权汇总，累计邻域特征从5D增加到25D时，ARI由0.3051提高至0.4418，NMI由0.4133提高至0.5982。然而，固定同为15D的A、B特征组合并未表现出一致的大幅优势；聚合层数、特征维数和输入块选择难以在这一组实验中完全拆开，不能简单归因为“更深的邻域一定更好”。

第二，聚类数设置改变了尺度比较的解释。E3c预设K网格显示ARI/NMI随K而变化，但这种扫描不是无标签模型选择；E3d原生AutoK在48个切片--尺度组合中完成无标签选择，其中10项触及冻结K=2至10搜索区间边界。事后独立评价表明，20D在4个累计尺度中拥有较高的12切片描述性ARI和NMI均值（0.4311与0.5942），而25D并未进一步提升到更高均值。不同尺度实际K可变，因此这些变化反映工作流程整体而非单一尺度机制的因果效应。

第三，空间边界评价呈现与整体聚类指标不同的趋势。E1/E2完整1200份Oracle-K源预测的事后边界分析显示，累计5D至25D的严格0-hop F1分别为0.1441至0.1179，而2-hop容差F1分别为0.2863至0.5333。由此可见，空间邻近容差会改变不同配置的比较结果，容差下较高的边界匹配率不能替代严格网格边界的准确定位。这是一项指标定义敏感性观察，不等于已通过实验确定特定的生物学“过度平滑”原因。

第四，与非空间GMM、Leiden及成熟空间方法BANKSY、BayesSpace比较时，CellCharter的表现取决于数据、K策略和指标，不能宣称对所有空间方法普遍占优。特别是本研究冻结的BayesSpace K-blind自动选K判据，在DLPFC的36项中有28项、乳腺癌的18项中有18项命中搜索上界，影响公平解释；不能为获得有利结果而在观察数据后改写原K规则。

第五，在六患者乳腺癌外部队列的K-blind评价中，CellCharter A15D相对5D基础配置的ARI宏平均差约为$+0.0458$，NMI差约为$+0.0583$，两项指标在6/6例患者上方向一致。但与BayesSpace比较时，ARI仅在4/6例患者较高，患者平均NMI差约为$-0.0139$。A15D在3/6例患者的ARI低于0.05，而在另一些患者表现明显较高，提示流程效果具有较强的样本异质性。此结果是将相同研究流程重新用于新组织队列所得的观察，不是将DLPFC已训练模型直接泛化到乳腺癌。

\subsection{研究局限}
本研究在解释层面主要受到以下因素限制：（1）DLPFC仅3位供体且存在供体内多个切片，乳腺癌仅6名患者，均不足以支持高精度总体推断；（2）DLPFC部分参考标注来源及旧开发样本的数据卫生问题需在论文数据来源中保留；（3）Oracle-K使用参考类别数，不能代表真正未知类别数的无标签应用；（4）部分实验中的邻域深度与输入维度同步变化，成熟基线的预处理及计算预算又存在方法间差异；（5）E3d的后验ARI/NMI以及E3b的1/2-hop边界分析是在已有模型结果之后追加的事后分析；（6）病理区域标签与皮层层级标签的语义不同，不能将跨组织ARI绝对值直接对比。

因此，本文展示的是基于已知模型的系统实证评价与可重复的研究流程，而非已验证的新核心聚类算法。在没有更多独立数据和精确因果控制前，不能对边界变化原因、参数最优性或临床应用价值作肯定性普遍结论。

\subsection{后续工作展望}
后续研究可在更多独立供体和组织类型上检验当前观察的适用性；在不读取评价标签的条件下比较更合适的聚类数选择策略；使用明确的等维和等计算预算对照，进一步区分尺度范围、特征维数和模型随机性；同时将组织层过渡区域的标注不确定性与边界评价结合，避免仅凭单一参考标签和单一指标判断算法表现。

以上属于对当前研究局限的未来研究方向，不表示本研究已经执行了上述扩展。本文当前工作已经完成相应冻结实验及事后分析，毕业论文写作应首先完整呈现已发生的研究事实，而不将计划工作当作已经取得的结果。
